
\section{Elementary elimination of surface base points}
\label{app:surface}

Elimination of indeterminacy by point blowups is a classical part of
the birational theory of smooth surfaces; see \cite{Bea96}. The next
lemma supplies exactly the resolution statement used in
\cref{lem:removal}. We include the argument to make explicit the finite
termination of the point-blowup process, using the colength of the local
base ideal as the decreasing positive integer.

\begin{lemma}\label[lemma]{lem:surface}
A rational map \(R:\PP^1\times\PP^1\dashrightarrow\PP^N\) is
holomorphic outside a finite set \(B\). A finite sequence of
point blowups over \(B\) gives a smooth projective surface
\(\mu:S'\to\PP^1\times\PP^1\) on which \(R\circ\mu\) extends
holomorphically.
The fibers of \(\mu\) are connected; its exceptional components
are smooth curves.
\end{lemma}

\begin{proof}
Clear denominators to represent \(R\) by a tuple of
bihomogeneous polynomials of the same bidegree, and remove
all common factors. The factors can be taken bihomogeneous:
the highest and lowest degrees in each variable group are
additive in products.

In any product affine chart, the dehomogenized tuple
\(h_0(x,y),\ldots,h_N(x,y)\) has no common nonconstant factor.
Otherwise bihomogenizing such a factor would give a common
factor of the original tuple. In \(\CC(x)[y]\), the tuple
therefore has gcd one. A univariate B\'ezout identity, with
denominators cleared, puts a nonzero polynomial in \(x\)
in the ideal \((h_0,\ldots,h_N)\). Interchanging the variables
gives a nonzero polynomial in \(y\). The common zero set
in the affine chart is finite. At each common zero, the ideal
in centered formal coordinates contains powers of both
coordinates and consequently has finite colength.
There are four product affine charts, so the global base set
\(B\) is finite.

If \(B=\varnothing\), the map is already holomorphic, and
\(S'=S\), \(\mu=\id_S\) give the conclusion. Assume from now on
that \(B\ne\varnothing\).

It remains to resolve such a base point. Let
\[
D=\CC[[x,y]],\qquad \fm=(x,y),
\]
and let \(I\subset D\) be the ideal of a representing holomorphic
tuple \(f_0,\ldots,f_N\) at the point. Assume \(I\) is proper
and has finite colength, and let
\(k=\min_\nu\operatorname{ord}(f_\nu)\ge1\).
Blow up the point. In a chart \((x,z)\mapsto(x,xz)\), divide
the pulled-back tuple by \(x^k\); the divided functions are
holomorphic. The other chart gives the same map, since the
two tuples differ by a unit on their overlap.

On the exceptional projective line, the remaining base points
are common zeros of the degree-\(k\) initial forms.
At least one such form is nonzero, so there are finitely many.
Fix one. A linear change of the old coordinates lets us place
it at \((x,z)=(0,0)\) in the displayed chart. Let
\[
J=(x^{-k}f_\nu(x,xz))_\nu\subset\CC[[x,z]],
\qquad D'=\CC[[x,z]]/J.
\]
A power \(x^\ell\) belongs to \(I\); applying the substitution
and division gives \(x^{\ell-k}\in J\).
Thus the class of \(x\) in \(D'\) is nilpotent.
Modulo \(x\), a nonzero degree-\(k\) initial form gives a
nonzero polynomial in \(z\) of degree at most \(k\). Its
germ at the new base point is \(z^q\) times a unit, with
\(1\le q\le k\). Therefore \(D'/(x)\) is spanned by
\(1,z,\ldots,z^{k-1}\).
Successively lifting modulo the powers of the nilpotent \(x\)
shows that \(D'\) is spanned by finitely many classes
\(x^az^b\), with \(0\le b<k\).

There is a well-defined \(\CC\)-linear map
\begin{equation}\label{eq:transform}
\fm^k/I\longrightarrow D',\qquad
h\longmapsto x^{-k}h(x,xz).
\end{equation}
Indeed a monomial \(x^ay^b\) of total degree at least \(k\)
maps to \(x^{a+b-k}z^b\). All exponents are nonnegative, and
the resulting orders tend to infinity with the original
orders, so the operation is defined on arbitrary formal series.
If \(h=\sum_\nu q_\nu f_\nu\) lies in \(I\), its transform is
\(\sum_\nu q_\nu(x,xz)x^{-k}f_\nu(x,xz)\in J\).
This proves descent to the quotient; no multiplicativity is
needed.

Each spanning monomial \(x^az^b\), \(b<k\), is the image of
\(x^{a+k-b}y^b\in\fm^k\), so \eqref{eq:transform} is surjective.
It follows that
\begin{equation}\label{eq:lengthdrop}
\dim_\CC D'
\le \dim_\CC(\fm^k/I)
=\dim_\CC(D/I)-\frac{k(k+1)}2
<\dim_\CC(D/I).
\end{equation}
In particular, the new base ideals have finite colength, and
that positive integer decreases strictly along every surviving
branch of the blowup process.

There are finitely many initial base points and finitely many
children at each step. The maximum initial colength bounds
the length of every branch by \eqref{eq:lengthdrop}. Thus the
entire tree is finite, and finitely many point blowups remove
all base points. The local divided tuples define the desired
holomorphic map.

Finally, a point blowup is proper, has a projective line over
its center, and is an isomorphism elsewhere.
The strict transform of a smooth curve remains smooth, and
a new exceptional line meets the strict transforms of the
components through the center. Induction gives smooth
exceptional components and connected fibers of the composite.
\end{proof}
